\section{Threats to Validity}
\label{sec:threat}

Validity depends on correct temporal and identity correspondences and on
interpreting the chronology within the reconstructed evidence. We
address these risks through controls in the reconstruction and separate
evaluation of its outputs.

\paragraph{Internal validity}
Because an incorrect impact pairing could propagate clock and identity
errors through the chronology, matching checks impulse compatibility,
footprint contact, available radar motion and consistency with
established offsets. Reconstruction uses recorder data and supplied road
context, with ground truth entering only at scoring. To reduce
confounding from configuration changes or different noise realizations,
the campaign fixes recordings, code and settings and reuses identical
per-recorder perturbations in corresponding conditions.

\paragraph{External validity}
Transfer to other sensing conditions depends on the structure of the
errors as well as their magnitude. The severe and extreme all-recorder
profiles apply sustained corruption to every radar throughout every
recording, while asymmetric runs assess partial degradation. Both quantify
sensitivity to the specified conditions without estimating their
deployment frequency. The error structure
matters for persistent radar ghosts, which could
accumulate support that independently redrawn clutter does not.
Confirmation, Doppler and cross-observer body checks constrain
unsupported hypotheses, although resilience to coherent ghosts
satisfying these checks remains unmeasured. The collision results are
obtained with impact reports available, while positioning and radar
evidence are perturbed.

\paragraph{Construct validity}
The oracle tests recovery of the defined manoeuvre and conflict onsets
by applying the same event rules to true trajectories. Criticality still
depends on the stated braking, response-delay and following-distance
assumptions. Since smoothing and manoeuvre confirmation use later
observations, semantic onsets are retrospective
estimates of those conditions. Relating them to recorded commands while
preserving source observations and identity provenance supports an
auditable interpretation, without equating a reconstructed conflict with
driver awareness or responsibility.

\paragraph{Conclusion validity}
Because the 1816 jobs repeat 22 recorded runs, seed-based intervals
describe perturbation variability within that set. Pooled differences
between single-recorder and all-recorder rows are interpreted
descriptively because their scenario weights differ. Local event scoring
uses ground-truth subject correspondence, whereas recorder links,
purity and merges separately assess physical identity. Likewise,
position errors are conditional on source-consistent states, with
separate false-entity and association counts exposing excluded errors.
The G0 geometric comparison pairs the same entity instants and measures
body reconstruction and fusion together. Reporting these dimensions
separately makes the reliability of the surrounding chronology visible
alongside collision identification.
